\section*{Chapter \ref{ch:s2:tail-hedging-and-long-volatility} --- Tail Hedging and Long Volatility}

\begin{solution}{exo:s2:tail-hedging-and-long-volatility:1}
$17 \times (-6.0) + 19.2 - 2.7 + 39.5 = -46.0\%$ of the portfolio, summed without compounding: the hedge lost money over the twenty years, and the
portfolio's growth fell from 3.0\% to 2.5\% a year while its worst drawdown fell from 73\% to 41\%.
\end{solution}

\begin{solution}{exo:s2:tail-hedging-and-long-volatility:2}
Each month's put is struck 5\% below that month's starting level, and each month the index falls about 4\%: no put ever ends in the money, and every
premium is lost while the index falls about 40\% over the year.
\end{solution}

\begin{solution}{exo:s2:tail-hedging-and-long-volatility:3}
Three times its cost: 3\% of notional.
\end{solution}

\begin{solution}{exo:s2:tail-hedging-and-long-volatility:4}
CLLZ's put spread protects only between 95\% and 97.5\% of the level; a fall of 30\% passes through the spread, and the short call gives up any
rebound. PPUT's single put protects all the way down.
\end{solution}

\begin{solution}{exo:s2:tail-hedging-and-long-volatility:5}
PPUT is a total-return index that reinvests dividends; the S\&P 500 price index does not. The true benchmark grows faster than the price index, so
the true shortfall of the hedged index is larger than the gap shown.
\end{solution}

\begin{solution}{exo:s2:tail-hedging-and-long-volatility:6}
A trend overlay turns short only after the trailing return has turned negative: it protects in long sell-offs and does nothing in a crash of days. A
put protects from its strike at once, whatever the path, and costs a premium whether or not a crash comes.
\end{solution}

\begin{solution}{exo:s2:tail-hedging-and-long-volatility:7}
Over that year the index loses 28.7\%, the static mix 21.0\% and the put programme 32.4\%, more than the unhedged index. Each month's put is struck
5\% below that month's start and the index falls less than that within most months, so the puts rarely pay and their premiums, dearer as vol rises,
add to the loss. It is Israelov's point in the chapter's model.
\end{solution}

\begin{solution}{exo:s2:tail-hedging-and-long-volatility:8}
Five calm years say nothing about a hedge built for rare crashes: on the synthetic market the programme bled in 17 years out of 20 and paid 19.2\% and
39.5\% in two of the others. Judge it by the portfolio over a history with crashes, and decide in advance whether its cost buys extra exposure
elsewhere.
\end{solution}

\begin{solution}{pb:s2:tail-hedging-and-long-volatility:1}
\begin{enumerate}
\item A position that pays in large, rare falls, bought at a cost; its bleed is that cost in calm periods, as a share of the portfolio a year.
\item The variance premium in the at-the-money vol and the skew premium of out-of-the-money puts.
\item Short-dated puts are the most reliable and most costly defence; bonds are unreliable; trend and quality did well in drawdowns.
\item Puts are often no better than statically holding less, unless timed right.
\item PPUT 7.7\% at 13.5\%, CLLZ 7.8\% at 14.6\%, the S\&P 500 price index 8.5\% at 18.4\% (without dividends).
\item 1987: $-10.6\%$, $-17.1\%$, $-21.8\%$; autumn 2008: $-6.7\%$, $-31.2\%$, $-30.1\%$; February--March 2020: $-0.4\%$, $-17.9\%$, $-19.9\%$.
\item The fall took twenty months; each month's put paid only for that month's fall beyond 5\%.
\item It protects only between its strikes.
\item Selling a hedge that has gained before expiry and replacing it at the new level.
\item From the table: 5\% monthly puts 2.5\% a year with a $-41\%$ drawdown; 20\% half-yearly 4.6\% and $-50\%$; monetised 2.2\% and $-62\%$; the
index 3.0\% and $-73\%$.
\item A 6.0\% bleed a year in 17 calm years; 19.2\%, $-2.7\%$ and 39.5\% in the crash years.
\item The replacement put, struck near the new lower spot, left the rest of the fall unprotected.
\item \textbf{Named result.} The monthly 5\% programme bleeds 6.0\% of the portfolio a year in calm years; the hedged portfolio grows at 2.5\% a year
against 3.0\% unhedged, with the worst drawdown cut from 73\% to 41\%.
\item Similar volatility (13.9\% against 13.2\%), a smaller drawdown ($-41\%$ against $-58\%$), lower growth (2.5\% against 2.8\%).
\item The trend overlay grows faster (6.3\%) with a deeper drawdown ($-56\%$).
\item Its crashes last ten days, the case puts are built for; slow bear markets are absent.
\item By the portfolio's growth rate over crashes of both kinds, and by what extra exposure its protection allows.
\item Fast and slow drawdowns, option prices on the smile, a static alternative, monetisation at the prices then.
\item Trend (or holding less).
\item Convexity in a crash, paid for in every other month.
\end{enumerate}
\end{solution}

\begin{solution}{iq:s2:tail-hedging-and-long-volatility:1}
Buyers want crash protection and sellers must be paid to provide it: index puts carry the variance premium and the skew premium, both largest for
out-of-the-money index puts.
\end{solution}

\begin{solution}{iq:s2:tail-hedging-and-long-volatility:2}
Match the two on volatility or on expected drawdown, then compare growth rates and drawdowns over a history with fast and slow falls, as Israelov did;
include costs and the timing of rolls.
\end{solution}

\begin{solution}{iq:s2:tail-hedging-and-long-volatility:3}
Follow the rule set in advance. Selling locks in the gain for reinvestment and resets protection near the new level; keeping the puts keeps the
protection against a further fall. A common compromise sells part and rolls the rest down.
\end{solution}

\begin{solution}{iq:s2:tail-hedging-and-long-volatility:4}
A fast crash, a slow year-long decline, a crash right after a roll and one right before, a volatility spike without a fall, and several calm years in a
row to test the budget.
\end{solution}

\begin{solution}{iq:s2:tail-hedging-and-long-volatility:5}
Each day: mark the puts on the surface, check the monetisation rule against the cost recorded at purchase, roll on schedule, size new puts to the
portfolio, record the bleed; with checks on stale option quotes and on roll dates.
\end{solution}

\begin{solution}{iq:s2:tail-hedging-and-long-volatility:6}
With $\log$ wealth growing by $\ln(1 + R)$ each period and $R$ log-normal with mean $\mu$ and variance $\sigma^2$, $E[\ln(1 + R)] \approx \mu -
\sigma^2/2$. A hedge that lowers $\mu$ by $c$ and $\sigma^2$ by more than $2c$ raises the growth rate: this is the case for a hedge that also allows a
larger position in the risky asset.
\end{solution}
